\subsection{n 阶导数}

设 \(\mathbf{f}\) 是实变向量函数，定义在区间 I 上、连续且可导。若导数 \(\mathbf{f}'\) 在一点 \(x_0 \in \mathrm{I}\) 的（相对于 I 的）一个邻域上存在，且它在点 \(x_0\) 可导，则其导数称为 \(\mathbf{f}\) 在点 \(x_0\) 的二阶导数，记作 \(\mathbf{f}''(x_0)\) 或 \(\mathrm{D}^2\mathbf{f}(x_0)\) 。若此二阶导数在 I 的每一点都存在（这蕴含 \(\mathbf{f}'\) 在 I 上存在且连续），则 \(x \mapsto \mathbf{f}''(x)\) 是一个向量函数，记作 \(\mathbf{f}''\) 或 \(\mathrm{D}^2\mathbf{f}\) 。我们同样地、递归地定义 \(n^{th}\) 阶导数（或 \(n\) 阶导数），即 \(\mathbf{f}\) 的相应导数，记作 \(\mathbf{f}^{(n)}\) 或 \(\mathrm{D}^n\mathbf{f}\) ；按定义，它在点 \(x_0 \in \mathrm{I}\) 的值是函数 \(\mathbf{f}^{(n-1)}\) 在点 \(x_0\) 的导数：这一定义预先假定诸导数 \(\mathbf{f}^{(k)}\) （\(k \leqslant n - 1\) 阶的）都在 \(x_0\) 相对于 I 的一个邻域上存在，且 \(\mathbf{f}^{(n-1)}\) 在点 \(x_0\) 可导。

我们将说 \(\mathbf{f}\) 是 \(n\) 次可导的（在一点 \(x_0\) 处；相应地，在一个区间中），如果它在这一点（相应地，在这个区间中）有 \(n^{th}\) 阶导数。则说 \(\mathbf{f}\) 在 I 上无穷次可导，如果对每个整数 \(n > 0\) 它在 I 上有 \(n\) 阶导数。

对 \(m\) 用归纳法可以看出

\[
\mathbf {D} ^ {m} (D ^ {n} \mathbf {f}) = \mathbf {D} ^ {m + n} \mathbf {f}.\tag{1}
\]

更精确地说，当 (1) 中的两项之一有定义时，另一项也有定义，且两者相等。

命题 1. 定义在区间 \(I \subset R\) 上、取值于给定的拓扑向量空间 E 中、且在 I 上有 \(n^{th}\) 阶导数的向量函数所成的集是 R 上的一个向量空间，且 \(f \mapsto D^{n}f\) 是从此空间到从 I 到 E 中的线性映射所成的向量空间中的线性映射。

我们证明公式

\[
\mathrm{D} ^ {n} (\mathbf {f} + \mathbf {g}) = \mathrm{D} ^ {n} \mathbf {f} + \mathrm{D} ^ {n} \mathbf {g}\tag{2}
\]

\[
\mathrm{D} ^ {n} (\mathbf {f} a) = \mathrm{D} ^ {n} \mathbf {f}. a\tag{3}
\]

当 f 和 g 在 I 上有 \(n^{th}\) 阶导数时（a 为常数），可对 n 用归纳法证明。

命题 2（“莱布尼茨公式”）. 设 E、F、G 是 R 上的三个拓扑向量空间，\((\mathbf{x}, \mathbf{y}) \mapsto [\mathbf{x}.\mathbf{y}]\) 是 \(E \times F\) 到 G 中的一个连续双线性映射。若 f（相应地，g）定义在区间 \(I \subset R\) 上、取值于 E（相应地，F）中且在 I 上有 \(n^{th}\) 阶导数，则 {[}f.g{]} 在 I 上有 \(n^{th}\) 阶导数，由公式

\[
\mathrm{D} ^ {n} [ \mathbf {f}. \mathbf {g} ] = [ \mathbf {f} ^ {(n)}. \mathbf {g} ] + \binom {n} {1} [ \mathbf {f} ^ {(n - 1)}. \mathbf {g} ^ {\prime} ] + \dots + \binom {n} {p} [ \mathbf {f} ^ {(n - p)}. \mathbf {g} ^ {(p)} ] + \dots + [ \mathbf {f}. \mathbf {g} ^ {(n)} ].\tag{4}
\]

给出。公式 (4) 可对 \(n\) 用归纳法证明（对二项式系数利用关系式 \(\binom{n}{p} = \binom{n-1}{p} + \binom{n-1}{p-1}\) ）。

同样可以验证下面的公式（其中的假设与命题 2 中的相同）：

\[
\begin{array}{c} [ \mathbf {f} ^ {(n)}. \mathbf {g} ] + (- 1) ^ {n - 1} [ \mathbf {f}. \mathbf {g} ^ {(n)} ] = D \big ([ \mathbf {f} ^ {(n - 1)}. \mathbf {g} ] - [ \mathbf {f} ^ {(n - 2)}. \mathbf {g} ^ {\prime} ] + \dots \\ + (- 1) ^ {n - 1} [ \mathbf {f}. \mathbf {g} ^ {(n - 1)} ] \big). \end{array}\tag{5}
\]

前面的命题是对在区间上 n 次可导的函数陈述的；我们把对在一点 n 次可导的函数的类似命题留给读者去陈述。

